%! Solving Quadratics by Factoring — Unit 3, Lesson 3.2 — Homework (Answer Key)
\documentclass[10pt]{article}
\usepackage{algebra2-article}
\usepackage{algebra2-key}   % pulls in -boxes; do NOT also load -boxes
\newcommand{\TallMath}[1]{$\displaystyle #1\rule[-1.4em]{0pt}{3.2em}$}
\newcommand{\lgrid}[4]{%
  \draw[step=1, linegray!40, very thin] (#1,#3) grid (#2,#4);
  \draw[->, semithick, charcoal] (#1,0)--(#2,0) node[below right, font=\scriptsize]{$x$};
  \draw[->, semithick, charcoal] (0,#3)--(0,#4) node[left, font=\scriptsize]{$y$};}

\begin{document}

\pageheader{Unit 3, Lesson 3.2}{Homework --- Answer Key}

\vspace{0.10in}

\begin{notesbox}{Practice}
Show your work: factored form first, then the roots. Move everything to one side before using the Zero
Product Property.

\begin{enumerate}[leftmargin=1.9em, itemsep=11pt]
  \item \textbf{Multiply, then factor (see the reverse).}
    \begin{center}
    \renewcommand{\arraystretch}{1.5}
    \begin{tabularx}{\linewidth}{@{}X X@{}}
    Multiply: $(x-3)(x+7)=$ \ans{$x^2+4x-21$} & Factor: $x^2+4x-21=$ \ans{$(x-3)(x+7)$} \\
    \end{tabularx}
    \end{center}
    What do you notice about the two answers? \ans{They undo each other --- multiplying builds the trinomial; factoring recovers the binomials.}

  \item \textbf{Factor completely.} (Watch for a GCF and the special patterns.)
    \begin{center}
    \renewcommand{\arraystretch}{1.6}
    \begin{tabularx}{\linewidth}{@{}X X X@{}}
    $x^2-3x-18=$ \ans{$(x-6)(x+3)$} & $x^2-36=$ \ans{$(x-6)(x+6)$} & $5x^2-20x=$ \ans{$5x(x-4)$} \\
    $x^2+8x+16=$ \ans{$(x+4)^2$} & $2x^2+5x-3=$ \ans{$(2x-1)(x+3)$} & $x^2+2x-15=$ \ans{$(x+5)(x-3)$} \\
    \end{tabularx}
    \end{center}

  \item \textbf{Solve by factoring.} List all roots.
    \begin{enumerate}[label=(\alph*), leftmargin=1.9em, itemsep=6pt]
      \item $x^2-7x+12=0$
        \begin{work}
          (x-3)(x-4) &= 0 \\
          x-3 = 0 \quad&\text{or}\quad x-4 = 0 \\
          x &= 3,\ 4
        \end{work}
      \item $x^2+5x=0$
        \begin{work}
          x(x+5) &= 0 \\
          x = 0 \quad&\text{or}\quad x+5 = 0 \\
          x &= 0,\ -5
        \end{work}
      \item $x^2=9x-20$
        \begin{work}
          x^2-9x+20 &= 0 \\
          (x-4)(x-5) &= 0 \\
          x-4 = 0 \quad&\text{or}\quad x-5 = 0 \\
          x &= 4,\ 5
        \end{work}
      \item $x^2-25=0$
        \begin{work}
          (x-5)(x+5) &= 0 \\
          x-5 = 0 \quad&\text{or}\quad x+5 = 0 \\
          x &= 5,\ -5
        \end{work}
    \end{enumerate}

  \item \textbf{Roots from a graph.} The graph of $f(x)=x^2-2x-8$ is shown.
    \begin{center}
    \begin{tikzpicture}[scale=0.5]
      \lgrid{-4}{5}{-10}{4}
      \begin{scope}\clip (-4,-10) rectangle (5,4);
        \draw[very thick, forest, domain=-2.6:4.6, samples=80, smooth] plot (\x,{(\x)*(\x) - 2*(\x) - 8});
      \end{scope}
      \fill[charcoal] (-2,0) circle (3pt);
      \fill[charcoal] (4,0) circle (3pt);
      \fill[charcoal] (0,-8) circle (3pt);
      \fill[charcoal] (1,-9) circle (3pt);
    \end{tikzpicture}
    \end{center}
    Factor $f(x)$: \ans{$(x+2)(x-4)$}. \; The roots are $x=$ \ans{$-2,\ 4$}, so the $x$-intercepts are
    \ans{$(-2,0),\ (4,0)$}. \; Does your algebra match the picture? \ans{yes}

  \item \textbf{Explain.} A classmate ``solves'' $x^2=5x$ by dividing both sides by $x$ to get $x=5$.
        What root did they lose, and what is the correct way? \ansline{They lost $x=0$. Correct: $x^2-5x=0 \Rightarrow x(x-5)=0 \Rightarrow x=0$ or $x=5$. Dividing by $x$ assumes $x\ne0$ and discards a real root.}
\end{enumerate}
\end{notesbox}

\begin{extensionbox}
\textbf{A bordered garden.} A square garden of side $x$ ft is surrounded so the whole plot measures
$(x+2)$ ft by $(x+5)$ ft. The total plot area is $70$ ft$^2$.
\begin{enumerate}[label=(\alph*), leftmargin=1.8em, itemsep=5pt]
  \item Write an equation for the plot's area, then rewrite it as a quadratic equal to $0$.
  \item Solve by factoring.
  \item One root cannot be a side length. State the garden's side length $x$ and explain why you rejected
        the other root.
\end{enumerate}
\ansline{(a) $(x+2)(x+5)=70 \Rightarrow x^2+7x+10=70 \Rightarrow x^2+7x-60=0$. \;\; (b) $(x+12)(x-5)=0 \Rightarrow x=-12$ or $x=5$. \;\; (c) Reject $x=-12$ (a side length can't be negative); the garden's side is $x=5$ ft. Check: $(7)(10)=70$ \checkmark.}
\end{extensionbox}

\begin{spiralbox}
\textbf{Coming next --- Lesson 3.3: Solving by Square Roots \& Completing the Square.} Factoring is fast
\emph{when it works}. But $x^2-6=0$ and $x^2+4x+1=0$ do not factor over the integers. Next you will solve
those by \emph{isolating a square} and by \emph{completing the square} --- turning any quadratic into a
perfect-square form you can undo.
\end{spiralbox}

\end{document}
